\documentclass[11pt,a4paper]{article}

% ── packages ──────────────────────────────────────────────
\usepackage[utf8]{inputenc}
\usepackage[T1]{fontenc}
\usepackage{amsmath,amssymb,amsthm,mathtools}
\usepackage{geometry}
\geometry{margin=2.5cm}
\usepackage{hyperref}
\hypersetup{colorlinks=true,linkcolor=blue!70!black,citecolor=blue!70!black,urlcolor=blue!70!black}
\usepackage{enumitem}
\usepackage{booktabs}
\usepackage{natbib}
\bibliographystyle{plainnat}
\usepackage{tikz}
\usetikzlibrary{arrows.meta,positioning,calc,decorations.pathmorphing}

% ── theorem environments ──────────────────────────────────
\newtheorem{definition}{Definition}
\newtheorem{proposition}{Proposition}
\newtheorem{remark}{Remark}

% ── macros ────────────────────────────────────────────────
\newcommand{\R}{\mathbb{R}}
\newcommand{\cC}{\mathcal{C}}
\newcommand{\cM}{\mathcal{M}}
\newcommand{\cL}{\mathcal{L}}
\newcommand{\fh}{\mathfrak{h}}
\newcommand{\fg}{\mathfrak{g}}
\DeclareMathOperator{\Lie}{Lie}
\DeclareMathOperator{\Ad}{Ad}
\DeclareMathOperator{\Pool}{Pool}
\DeclareMathOperator{\diag}{diag}

% ══════════════════════════════════════════════════════════
\title{\textbf{Compositional Dynamic Depth in Neural Networks}\\[4pt]
\large A Symmetry-Theoretic Foundation and Experimental Programme}
\author{Giulio Ruffini\\
\small Neuroelectrics, Starlab -- Barcelona, Spain\\
\small\texttt{giulio.ruffini@neuroelectrics.com}}
\date{Technical Note -- Draft, March 2026}

\begin{document}
\maketitle

% ──────────────────────────────────────────────────────────
\begin{abstract}
Standard deep networks apply every layer uniformly to every input.
We argue that \emph{depth should be dynamic and input-dependent},
controlled by the compositional complexity of the data-generating process.
Building on the Lie-pseudogroup framework of compositional symmetry
\citep{Ruffini2025entropy,Ruffini2025compositional},
we show that the hierarchical flag
$G=H_0\supset H_1\supset\cdots\supset H_L$
of a generative model naturally prescribes which layers
an agent needs to engage for a given input:
layers corresponding to inactive symmetry scales should be gated off.
We formulate a \emph{symmetry-gated residual network} in which
multiplicative gates have precise geometric meaning---they measure the
presence of unresolved structure at each scale of the pseudogroup
hierarchy.
We then lay out a concrete experimental programme using
Blender-generated datasets with known compositional structure
to validate the theory.
\end{abstract}

% ══════════════════════════════════════════════════════════
\section{Motivation: Why Dynamic Depth?}\label{sec:motivation}

In a standard $L$-layer residual network the update rule is
\begin{equation}\label{eq:resnet}
  h_{\ell+1} = h_\ell + \Delta_\ell(h_\ell),
  \qquad \ell = 0,1,\ldots,L-1,
\end{equation}
where $\Delta_\ell$ is the residual block at layer $\ell$.
Every input traverses every layer with equal weight.
Yet the compositional complexity of natural data varies enormously:
a static scene with fixed illumination exercises far fewer
generative degrees of freedom than an articulated body
under changing viewpoint and lighting.

Several lines of work address this mismatch:
Adaptive Computation Time \citep{Graves2016},
Universal Transformers \citep{Dehghani2019},
Mixture-of-Depths \citep{Raposo2024},
and LayerSkip \citep{Elhoushi2024}.
All share the engineering goal of saving compute by skipping layers.
What is missing is a \emph{principled theory that predicts
when and why depth should be reduced}---one that connects
the gate pattern to the structure of the data-generating process.

This note proposes that the Lie-pseudogroup framework of
compositional symmetry
\citep{Ruffini2025entropy,Ruffini2025compositional}
provides exactly such a theory.
The central claim is:

\medskip
\noindent\fbox{\parbox{0.95\linewidth}{%
\textbf{Thesis.}\quad
The effective depth a network needs for a given input
is determined by the number of active levels in the
compositional factorisation of the data-generating symmetry group.
Layers aligned with inactive symmetry scales carry near-zero residuals
and should be gated off.
Dynamic depth is not an architectural trick---it is a
consequence of compositional structure in the world.}}
\medskip

% ══════════════════════════════════════════════════════════
\section{Theoretical Framework}\label{sec:theory}

We briefly recall the essential apparatus;
full details are in \citet{Ruffini2025compositional}
and \citet{Ruffini2025entropy}.

\subsection{Lie-Pseudogroup Generative Models}

\begin{definition}[Generative model]
A generative model is a smooth map
$f:\cC\to\R^X$ from an $M$-dimensional configuration manifold
($M\ll X$) to observations.
It is a \emph{Lie generative model} if a Lie pseudogroup $G$
acts on $\cC$ and $\R^X$ such that
$f(c) = f(\gamma\cdot c_0) = \gamma\cdot I_0$
for $\gamma\in G$ and a reference observation $I_0$.
\end{definition}

Near the identity, elements of $G$ factor as
\begin{equation}\label{eq:factor}
  \gamma = \gamma^{(L)}\gamma^{(L-1)}\cdots\gamma^{(1)},
  \qquad \gamma^{(k)}\in A_k,
\end{equation}
where each $A_k$ is a local sub-pseudogroup generated by
``level-$k$'' infinitesimal symmetries.
This is the formal expression of compositional nesting:
complex deformations are built by composing a few primitive
moves close to the identity.

\subsection{The Symmetry Flag and Reduced Manifolds}

Choosing the factorisation order in \eqref{eq:factor}
induces a \emph{flag of sub-pseudogroups},
\begin{equation}\label{eq:flag}
  G = H_0 \supset H_1 \supset \cdots \supset H_L,
  \qquad H_k := A_{k+1}\cdots A_L,
\end{equation}
obtained by successively omitting generators.
Each step $H_{k-1}\to H_k$ declares a new set of symmetry
directions ``already explained'' at coarser scale.
The quotients $H_{k-1}/H_k \simeq A_k$ collect the
\emph{residual directions} that remain to be explained at level~$k$.

Passing to orbit spaces produces nested \emph{reduced manifolds}:
\begin{equation}\label{eq:nested}
  \cM_0 \supset \cM_1 := \cM_0/H_1 \supset \cdots
  \supset \cM_L := \cM_{L-1}/H_L,
  \qquad \dim\cM_{k+1} < \dim\cM_k.
\end{equation}
Motion within $\cM_k$ is driven only by
residual generators in~$H_k$.

\subsection{Dynamic Depth from the Flag}\label{sec:dynamic}

Now consider an input stream $I_{\theta(t)}$ in which only
generators at levels $1,\ldots,K$ are active
(i.e.\ $\gamma^{(k)}\approx e$ for $k>K$).
Then by construction:

\begin{itemize}[nosep]
\item The data lie on $\cM_K$, the reduced manifold obtained
  by quotienting out the inactive levels.
\item At any level $k>K$, the residual
  $\varepsilon_k\in H_{k-1}/H_k$ is near the identity:
  there is no unexplained structure at that scale.
\item A world-tracking agent that mirrors this hierarchy
  should produce near-zero updates at layers $k>K$.
\end{itemize}

This is the geometric origin of dynamic depth:
\textbf{the effective depth equals the number of active
compositional levels in the generative process.}

% ══════════════════════════════════════════════════════════
\section{Symmetry-Gated Residual Architecture}\label{sec:arch}

We now translate the theory into a concrete neural architecture.

\subsection{General Formulation}

Replace the standard residual rule \eqref{eq:resnet} with
\begin{equation}\label{eq:gated}
  \boxed{
  h_{\ell+1} = h_\ell
    + \alpha_\ell(x,\, h_\ell)\;\Delta_\ell(h_\ell),
  \qquad \alpha_\ell \in [0,1].
  }
\end{equation}
When $\alpha_\ell\approx 0$ the layer is effectively skipped;
when $\alpha_\ell\approx 1$ it contributes fully.
This is a soft, differentiable version of dynamic depth.

\paragraph{Geometric interpretation.}
In the symmetry-theoretic picture, layer $\ell$ corresponds
to level $k(\ell)$ of the flag.
The gate $\alpha_\ell$ measures whether the input has
unresolved structure in the quotient $H_{k-1}/H_k$:
\begin{equation}
  \alpha_\ell \;\approx\; \|\text{residual component in }
  \Lie(A_{k(\ell)})\|.
\end{equation}
This is not a free scalar learned by gradient descent alone---it
has a \emph{predicted} relationship to the compositional complexity
of the input.

\subsection{Gate Design Choices}

\paragraph{Granularity.}
The gate can operate at different granularities:
\begin{center}
\begin{tabular}{@{}lll@{}}
\toprule
Granularity & Gate shape & Interpretation \\
\midrule
Per-sequence & $\alpha_\ell(x)\in\R$ & One depth profile per input \\
Per-token & $\alpha_{\ell,t}(h_{\ell,t})\in\R$ & Token-wise depth \\
Per-channel/head & $\alpha_{\ell,t,c}\in\R^{d}$ & Fine-grained routing \\
\bottomrule
\end{tabular}
\end{center}
For the Blender experiments (Section~\ref{sec:experiments}),
per-sequence gating suffices because the generative hierarchy
is a property of the scene, not of individual pixels.
For language, per-token gating is more natural.

\paragraph{Soft vs.\ hard gates.}
Soft gates (sigmoid) are easy to train and allow fractional
contribution. Hard gates (top-$k$, straight-through) save
FLOPs but require tricks for gradient flow.
We recommend starting with soft gates and adding a
\emph{compute penalty} to encourage sparsity:
\begin{equation}\label{eq:loss}
  \cL = \cL_{\text{task}}
  + \lambda\sum_{\ell=1}^{L}\alpha_\ell,
\end{equation}
where $\lambda>0$ controls the depth--accuracy trade-off.

\paragraph{Budgeted variant.}
For a fixed compute budget $C$, one can impose
$\sum_\ell\alpha_\ell = C$
using a softmax or Sinkhorn normalisation over layers.
This is closer in spirit to Mixture-of-Depths.

\subsection{Concrete Gate Module}\label{sec:gate-module}

A minimal implementation:
\begin{equation}\label{eq:gate-impl}
  \alpha_\ell = \sigma\!\bigl(
    w_\ell^\top \Pool(h_\ell) + b_\ell
  \bigr),
\end{equation}
where $\Pool$ is global average pooling (for images) or
\texttt{[CLS]} extraction (for transformers),
$w_\ell\in\R^d$, $b_\ell\in\R$, and $\sigma$ is the sigmoid.
This adds $d+1$ parameters per layer---negligible overhead.

For token-wise gating:
\begin{equation}
  \alpha_{\ell,t} = \sigma(W_\ell\, h_{\ell,t} + b_\ell),
  \qquad W_\ell\in\R^{1\times d}.
\end{equation}

\subsection{Relationship to Predictive Coding}

The predictive hierarchy of \citet{Ruffini2025compositional}
prescribes a richer per-layer algorithm than a simple gate:
\begin{equation}\label{eq:predcoding}
  \text{predict} \;\to\;
  \text{compare} \;\to\;
  \text{update locally via } P_k \;\to\;
  \text{canonicalise} \;\to\;
  \text{coarse-grain via } \mathcal{C}_{k\to k+1} \;\to\;
  \text{pass up}.
\end{equation}
Here $P_k:\fg\to\fh_k$ is the orthogonal projector onto
$\Lie(H_k)$, and $Q_k = I-P_k$ extracts the quotient directions.
Each level absorbs what it can explain
($P_k\eta_k$) and forwards only the unresolved remainder
($Q_k\eta_k$) after canonicalisation and coarse-graining.

The simple gate $\alpha_\ell$ in \eqref{eq:gated} is a
\emph{scalar summary} of this richer process:
$\alpha_\ell\approx \|Q_{k(\ell)}\eta_{k(\ell)}\|$.
Building architectures that implement the full
predict--compare--canonicalise--coarse-grain loop
at each level is a natural next step
(see Section~\ref{sec:future}).


% ══════════════════════════════════════════════════════════
\section{Experimental Programme}\label{sec:experiments}

The key advantage of this framework is that we can
\emph{predict} gate patterns from the known generative hierarchy,
and then test those predictions empirically.
This requires a dataset where the compositional structure
is fully controlled.

\subsection{The Blender Cat Rig as a Testbed}\label{sec:blender}

Following Appendix~C of \citet{Ruffini2025compositional},
we use a production-quality Blender rig
(e.g.\ the cat from the \emph{Flow} project,
or any rigged character) whose configuration space decomposes as
\begin{equation}
  \cC \;\simeq\;
  \cC_{\text{camera}} \times
  \cC_{\text{body}} \times
  \cC_{\text{spine}} \times
  \cC_{\text{limbs}} \times
  \cC_{\text{face}} \times
  \cC_{\text{fur}} \times
  \cC_{\text{light}} \times
  \cC_{\text{env}}.
\end{equation}
The eight levels correspond to sub-pseudogroups $A_1,\ldots,A_8$
(Table~\ref{tab:levels}).

\begin{table}[h]
\centering
\caption{Lie-pseudogroup hierarchy for the Blender cat rig.}\label{tab:levels}
\smallskip
\begin{tabular}{@{}clll@{}}
\toprule
Level $k$ & Generators $A_k$ & Blender construct & $\dim A_k$ \\
\midrule
1 & Camera $SE(3)$ + lens & \texttt{camera.matrix\_world}, focal length & 7 \\
2 & Global body/root $SE(3)$ & \texttt{pose.bones["root"]} & 6 \\
3 & Torso/spine chain (PoE) & Spine bones, FK/IK & $\sim$12 \\
4 & Limbs/paws/tail (PoE) & Limb bones, IK solvers & $\sim$24 \\
5 & Facial morphology & Shape keys (blendshapes) & $\sim$20 \\
6 & Fur/appearance & Material nodes, groom PCs & $\sim$6 \\
7 & Illumination & Light transforms, SH coeffs & $\sim$12 \\
8 & Environment/camera jitter & World rotation, exposure & $\sim$4 \\
\bottomrule
\end{tabular}
\end{table}

\subsection{Dataset Generation Pipeline}\label{sec:datagen}

We propose generating the dataset entirely through
Blender's Python API (\texttt{bpy}), which avoids
any need to ``patch'' Blender and can be run headless
on a compute cluster.

\paragraph{Step 1: Acquire or build a rigged character.}
Options in decreasing order of realism:
\begin{enumerate}[nosep]
\item Download a production-ready rigged cat
  (e.g.\ from BlenderKit, Sketchfab, or TurboSquid---many
  are free under CC licences).
  Look for rigs with armature, shape keys, and
  shader-based materials.
\item Use the free ``Cat'' asset from the Blender demo files.
\item Build a simpler parametric model (e.g.\ an articulated
  stick figure) for a first proof-of-concept.
\end{enumerate}

\paragraph{Step 2: Python script for controlled variation.}
A Blender Python script that:
\begin{enumerate}[nosep]
\item Defines the $A_k$ parameter ranges
  (joint angle limits, shape-key ranges, material parameter
  ranges, light position/colour ranges, camera orbit).
\item For each sample, draws $\theta^{(k)}\sim\text{Uniform}(A_k)$
  for the \emph{active} levels and holds inactive levels at
  their default (rest) values.
\item Writes pose bone transforms via
  \texttt{pose.bones[name].rotation\_quaternion},
  shape key values via \texttt{key\_blocks[name].value},
  material parameters via node socket values,
  light transforms via \texttt{light.matrix\_world},
  and camera via \texttt{camera.matrix\_world}
  and \texttt{camera.data.lens}.
\item Renders to an image buffer using Cycles or Eevee
  and saves the image together with a metadata JSON
  recording all parameter values and
  which levels were active.
\end{enumerate}

\paragraph{Step 3: Construct controlled datasets.}
For each subset $S\subseteq\{1,\ldots,8\}$ of active levels,
generate $N$ images with only those levels varying.
The key conditions are:

\begin{center}
\begin{tabular}{@{}ll@{}}
\toprule
Condition name & Active levels \\
\midrule
\textsc{Static} & $\emptyset$ (single image, repeated) \\
\textsc{CameraOnly} & $\{1\}$ \\
\textsc{CameraBody} & $\{1,2\}$ \\
\textsc{CameraBodySpine} & $\{1,2,3\}$ \\
\textsc{FullPose} & $\{1,2,3,4\}$ \\
\textsc{PoseFace} & $\{1,2,3,4,5\}$ \\
\textsc{AllGeometric} & $\{1,\ldots,5\}$ \\
\textsc{AllPlusAppearance} & $\{1,\ldots,7\}$ \\
\textsc{Everything} & $\{1,\ldots,8\}$ \\
\bottomrule
\end{tabular}
\end{center}

\paragraph{Step 4: Run headless.}
All rendering can be launched from the command line:
\begin{verbatim}
  blender --background scene.blend \
          --python generate_dataset.py -- \
          --condition CameraBody --n_samples 5000
\end{verbatim}
No GUI needed; this runs on any Linux node with a GPU.

\subsection{Training Protocol}\label{sec:training}

\paragraph{Architecture.}
Take a standard ResNet-50 (or Vision Transformer)
and insert a gate module \eqref{eq:gate-impl}
at every residual block.
Train on the \textsc{Everything} condition
with the gated loss~\eqref{eq:loss}.

\paragraph{Baselines.}
\begin{enumerate}[nosep]
\item \emph{Standard ResNet-50}: no gates (fixed depth).
\item \emph{Random gates}: $\alpha_\ell$ sampled uniformly,
  to control for the regularisation effect of the penalty.
\item \emph{ACT-style halting}: a single halting probability
  rather than per-layer gates.
\end{enumerate}

\subsection{Measurements and Predictions}\label{sec:predictions}

The theory makes the following \textbf{testable predictions}:

\begin{enumerate}
\item \textbf{Gate--complexity alignment.}\quad
  When the trained model is evaluated on the controlled conditions,
  the gate vector $(\alpha_1,\ldots,\alpha_L)$ should exhibit a
  clear transition: gates corresponding to layers
  aligned with inactive levels should be significantly
  smaller than those aligned with active levels.

  \emph{Metric:} For each condition with $K$ active levels,
  compute the mean gate activation for ``active'' vs.\
  ``inactive'' layer groups.
  The theory predicts a monotone relationship:
  more active levels $\Rightarrow$ more layers engaged.

\item \textbf{Depth--complexity curve.}\quad
  Define the \emph{effective depth}
  $D_{\text{eff}}(x) := \sum_\ell \alpha_\ell(x)$.
  The theory predicts $D_{\text{eff}}$ increases with the
  number of active generative levels:
  \begin{equation}
    D_{\text{eff}}(\textsc{CameraOnly})
    < D_{\text{eff}}(\textsc{CameraBody})
    < \cdots
    < D_{\text{eff}}(\textsc{Everything}).
  \end{equation}

\item \textbf{Residual structure.}\quad
  The residual norms $\|\Delta_\ell(h_\ell)\|$ at
  gated-off layers should be small not merely because the gate
  suppresses them, but because the \emph{pre-gate} residual
  itself carries little energy at that scale.
  This distinguishes our account from pure regularisation artifacts.

\item \textbf{Generalisation to novel compositions.}\quad
  Train on conditions up to \textsc{FullPose} ($K=4$),
  then evaluate on \textsc{PoseFace} ($K=5$) without retraining.
  The theory predicts that gate~5 (face) will
  activate to accommodate the new level,
  demonstrating compositional generalisation of the depth profile.
\end{enumerate}


% ══════════════════════════════════════════════════════════
\section{Relationship to Prior Work}\label{sec:related}

\paragraph{Adaptive Computation Time (ACT) \citep{Graves2016}.}
ACT learns a scalar halting probability per step in a
recurrent (shared-weight) architecture.
Our framework differs in three ways:
(i) layers are \emph{distinct} (different $A_k$ at each level),
so dynamic depth means selecting which \emph{compositions of
different functions} to apply;
(ii) gates have geometric meaning tied to the symmetry hierarchy;
(iii) the prediction of gate patterns is derived from the
generative model, not learned from scratch.

\paragraph{Mixture-of-Depths \citep{Raposo2024}.}
MoD routes a subset of tokens through each layer via top-$k$
selection. This is per-token, hard-gated depth allocation.
Our soft per-layer gating is complementary;
the budgeted variant of \eqref{eq:loss} with a softmax
over layers converges to a similar mechanism.

\paragraph{Neural ODEs \citep{Chen2018}.}
Neural ODEs make depth continuous; adaptive step-size solvers
implement a form of dynamic depth.
Our gate magnitude $\alpha_\ell$ plays the role of a
step size in a discretised ODE.
The connection to ODE stability analysis
(Lyapunov functions, ISS bounds) is made explicit in
\citet{Ruffini2025compositional}, Equation~(5).

\paragraph{Deep Equilibrium Models \citep{Bai2019}.}
DEQs take dynamic depth to the limit of infinite (shared) layers
and find a fixed point.
Our framework suggests that the fixed-point view
is appropriate when the compositional hierarchy is
\emph{recursive} (shared block), but
\emph{compositional} hierarchies with distinct levels
require distinct layers---making gated residual networks
the natural architecture.


% ══════════════════════════════════════════════════════════
\section{Future Directions}\label{sec:future}

\begin{enumerate}
\item \textbf{Full predictive-coding architecture.}\quad
  Replace the scalar gate with the full
  predict--compare--project--canonicalise--coarse-grain
  loop of \eqref{eq:predcoding}.
  Each layer would contain:
  (a) a predictor $\hat{I}_k = \hat\gamma_k\cdot I_0$,
  (b) a comparator computing $\delta I_k = I - \hat{I}_k$,
  (c) a linear residual fit
  $\eta_k \approx \sum_a \eta_k^a V_a^{(k)}$,
  (d) a projection $P_k\eta_k$ (local update)
  and $Q_k\eta_k$ (upward message),
  (e) a coarse-graining operator $\mathcal{C}_{k\to k+1}$.
  This is architecturally richer than a gated ResNet and
  directly implements the pseudogroup hierarchy.

\item \textbf{Symmetry discovery.}\quad
  Rather than prescribing the flag,
  learn the generators $\{T_k\}$ and the factorisation
  from data using methods such as LieGG \citep{Moskalev2022}
  or the approach of \citet{Hu2024}.
  The learned flag would then \emph{predict} the gate pattern.

\item \textbf{Language and transformers.}\quad
  Syntactic/semantic compositionality in language provides
  a natural (if less controlled) hierarchy.
  Per-token gating in a transformer could reveal whether
  depth allocation correlates with measures of
  linguistic compositional complexity
  (syntactic depth, number of variable bindings,
  logical nesting).

\item \textbf{Brain signals.}\quad
  EEG/neural signals exhibit varying temporal complexity.
  A dynamic-depth model for neural decoding where the model
  allocates more compositional processing to complex signal
  epochs would connect to both the neuroscience of
  processing difficulty and to KT's account of
  structured experience.
\end{enumerate}


% ══════════════════════════════════════════════════════════
\section{Summary}

The standard dynamic-depth literature (ACT, MoD, LayerSkip)
approaches the problem from an engineering direction:
save compute by skipping layers.
We approach it from the opposite direction:
the Lie-pseudogroup framework of compositional symmetry
\emph{predicts} when and why depth should be dynamic.

The experimental programme proposed here---using Blender rigs
with fully controlled compositional structure---provides a
direct test: train a gated-depth network, then verify that
the learned gate pattern matches the symmetry-predicted
depth profile.
If it does, we have the first empirical demonstration that
dynamic depth in neural networks reflects the compositional
structure of the data-generating process.

The contribution is not another gated-depth mechanism.
It is a \emph{theory that explains when and why dynamic depth
emerges}, grounded in the geometry of Lie pseudogroups,
together with a concrete plan to test it.


% ══════════════════════════════════════════════════════════
\begin{thebibliography}{20}
\bibitem[Bai et~al.(2019)]{Bai2019}
S.~Bai, J.~Z. Kolter, and V.~Koltun.
\newblock Deep equilibrium models.
\newblock In \emph{NeurIPS}, 2019.

\bibitem[Chen et~al.(2018)]{Chen2018}
R.~T.~Q. Chen, Y.~Rubanova, J.~Bettencourt, and D.~Duvenaud.
\newblock Neural ordinary differential equations.
\newblock In \emph{NeurIPS}, 2018.

\bibitem[Dehghani et~al.(2019)]{Dehghani2019}
M.~Dehghani, S.~Gouws, O.~Vinyals, J.~Uszkoreit, and \L.~Kaiser.
\newblock Universal transformers.
\newblock In \emph{ICLR}, 2019.

\bibitem[Elhoushi et~al.(2024)]{Elhoushi2024}
M.~Elhoushi, A.~Shrivastava, D.~Liskovich, B.~Hosseini,
  R.~Rabbat, and C.-J.~Hsieh.
\newblock {LayerSkip}: Enabling early exit inference and
  self-speculative decoding.
\newblock \emph{arXiv:2404.16710}, 2024.

\bibitem[Graves(2016)]{Graves2016}
A.~Graves.
\newblock Adaptive computation time for recurrent neural networks.
\newblock \emph{arXiv:1603.08983}, 2016.

\bibitem[Hu et~al.(2024)]{Hu2024}
L.~Hu, Y.~Li, and Z.~Lin.
\newblock Symmetry discovery for different data types.
\newblock \emph{arXiv:2410.09841}, 2024.

\bibitem[Moskalev et~al.(2022)]{Moskalev2022}
A.~Moskalev, A.~Sepliarskaia, I.~Sosnovik, and A.~Smeulders.
\newblock {LieGG}: Studying learned {L}ie group generators.
\newblock \emph{arXiv}, 2022.

\bibitem[Raposo et~al.(2024)]{Raposo2024}
D.~Raposo, S.~Ritter, B.~Richards, T.~Lillicrap,
  P.~Conway~Humphreys, and A.~Santoro.
\newblock Mixture-of-depths: Dynamically allocating compute
  in transformer-based language models.
\newblock \emph{arXiv:2404.02258}, 2024.

\bibitem[Ruffini(2025a)]{Ruffini2025compositional}
G.~Ruffini.
\newblock Compositional symmetry as compression:
  {L}ie-pseudogroup structure in algorithmic agents.
\newblock \emph{arXiv:2510.10586}, 2025.

\bibitem[Ruffini et~al.(2025b)]{Ruffini2025entropy}
G.~Ruffini, F.~Castaldo, and J.~Vohryzek.
\newblock Structured dynamics in the algorithmic agent.
\newblock \emph{Entropy}, 27(1):90, 2025.
\newblock \doi{10.3390/e27010090}.

\end{thebibliography}

\end{document}
